\section{Dataset Construction}\label{app:dataset}

This appendix details the construction of the labeled evaluation dataset used
throughout the paper.

%----------------------------------------------------------------------
\subsection{Agent Address Extraction}

A total of 1{,}505 ERC-8004-registered agent addresses were extracted from the
Base chain agent registry contract. These addresses serve as ground-truth
positives: each is provably an autonomous agent by virtue of its on-chain
registration under the ERC-8004 standard.

%----------------------------------------------------------------------
\subsection{Transaction Collection}

USDC transfer events involving the extracted agent addresses were queried via
Dune Analytics. Due to API rate limits, queries were executed in 16 batches of
100 addresses each. The raw result comprised 81{,}904 transaction rows,
subsequently preprocessed and stored in Apache Parquet format for efficient
downstream processing.

%----------------------------------------------------------------------
\subsection{Labeling Strategy}

\begin{itemize}[leftmargin=*]
  \item \textbf{Positive class (agent).}
    Any address present in the ERC-8004 registry is labeled as an agent.
    This constitutes a deterministic ground-truth label---no probabilistic
    inference is required.
  \item \textbf{Negative class (human counterparty).}
    Counterparty addresses that are \emph{not} in the ERC-8004 registry and
    that appear in $\geq 5$ transactions with known agents are labeled as
    human with a confidence of 0.7. The transaction-count threshold filters
    out ephemeral or bot-like addresses that interact only once, while the
    sub-unit confidence reflects residual uncertainty (some counterparties
    may themselves be unregistered agents).
\end{itemize}

%----------------------------------------------------------------------
\subsection{Final Evaluation Set}

After filtering and deduplication:
\begin{itemize}[leftmargin=*]
  \item 665 agent addresses retained in the final evaluation set;
  \item 1{,}069 cleaned human counterparty addresses.
\end{itemize}

The reduction from 1{,}505 to 665 agents is due to addresses with insufficient
transaction history (fewer than the minimum required for signal computation).

%----------------------------------------------------------------------
\subsection{Attack Injection}

To evaluate adversarial robustness, synthetic attack transactions were injected
into the dataset:
\begin{itemize}[leftmargin=*]
  \item 6{,}050 synthetic transactions generated;
  \item 82 unique attack addresses, all prefixed with \texttt{0xdead} to
    guarantee no collision with real on-chain addresses;
  \item Injection rate: 6.07\% of the combined dataset;
  \item Attack scenarios include wash trading rings, Sybil clusters, and
    layering chains (see \Cref{sec:evaluation} for results).
\end{itemize}

%----------------------------------------------------------------------
\subsection{Data Files}

The released dataset consists of three Parquet files:

\begin{table}[h]
\centering
\caption{Dataset files and their contents.}
\label{tab:data-files}
\begin{tabular}{lp{8cm}}
\toprule
\textbf{File} & \textbf{Description} \\
\midrule
\texttt{transactions\_dune.parquet} &
  81{,}904 raw USDC transfer events (sender, receiver, value, timestamp,
  block number). \\
\texttt{labels\_dune.parquet} &
  Ground-truth labels for 665 agents and 1{,}069 human counterparties
  with associated confidence scores. \\
\texttt{attack\_injection\_dataset.parquet} &
  6{,}050 synthetic attack transactions with 82 attack addresses and
  scenario annotations. \\
\bottomrule
\end{tabular}
\end{table}
